\subsection{Schwächen und Grenzen}
\label{subsec:schwaechen-grenzen}

Für die Einordnung sind drei Fälle zu unterscheiden. Eine \emph{technische
Schwäche} ist ein Nachteil der vorhandenen Lösung, eine \emph{Scope-Grenze}
eine bewusste Auslassung und ein \emph{nicht verifizierter Bereich} eine
fehlende Nachweisabdeckung. Erstere wirkt unmittelbar auf die Bewertung;
letztere beiden sind nicht automatisch Implementierungsfehler.

Technisch erhöht die Bündelung im Master-Repository dessen interne
Komplexität. Änderungen an gemeinsamen Pfaden, Generator und Feature-Katalog
können mehrere Profile betreffen und verlangen Regressionstests über die
Varianten. Der derzeit kleine Katalog begrenzt dieses Risiko, beseitigt es aber
nicht. Hinzu kommen zwei konkrete Architekturgrenzen: Die gemeinsamen
JSON-Schemata werden von Frontend und Backend noch nicht zur Laufzeit
konsumiert, sodass Schnittstellenänderungen manuell synchronisiert werden
müssen (Abschnitt~\ref{subsec:shared-contracts-schnittstellen}). Zudem trennt
die direkte CLI-Auswahl \texttt{optional} und \texttt{selectable} weniger
streng als der interaktive Dialog (Abschnitt~\ref{subsec:optionale-capabilities}).
Beide Punkte erzeugen Wartungs- beziehungsweise Fehlbedienungsrisiken, ohne
die geprüften Presets grundsätzlich infrage zu stellen. Hinzu kommt ein
Prozessbefund: Auf dem untersuchten HEAD bestanden die Webprofile in der
externen Matrix, während der Tooling-Job, die Desktop-Profiltests und der
Windows-Installationsschritt fehlschlugen; einem lokal bestandenen
Tooling-Kontrolllauf steht damit ein widersprüchlicher Remote-Nachweis
gegenüber (Abschnitt~\ref{subsec:continuous-integration}).

Die Profilgrenzen sind dagegen beabsichtigt: \texttt{web-only} enthält kein
Backend, \texttt{desktop-local} keine Cloud-API, und PostgreSQL ist nur mit
Backendprofilen vereinbar. Ebenso gehören Fachlogik, generische
Authentifizierung oder Rollenmodelle, eine verpflichtende Datenbank, ein
bestimmter Cloud-Provider sowie produktspezifische native Kommandos nicht zur
Baseline. Daraus entsteht Anpassungsaufwand für konkrete Produkte, aber keine
Abweichung von der Abgrenzung in Abschnitt~\ref{subsec:abgrenzung}.

Nicht abschließend verifiziert sind die in
Abschnitt~\ref{subsec:extern-verbleibende-pruefungen} vorgesehenen externen
Punkte. Die aktuellen CI-Läufe belegen Compose-Validierung und beide
Master-Image-Builds, nicht jedoch den Start des Compose-Verbunds oder eine
produktive Bereitstellung. Offen bleiben ferner Release Validation und Branch
Protection. Linux- und macOS-Pakete wurden auf dem untersuchten HEAD
unsigniert gebaut; der Windows-Job scheiterte. Signing, macOS-Notarisierung,
Veröffentlichung und produktive Cloud-Bereitstellung bleiben bewusst produkt-
oder infrastrukturspezifisch
\parencite{kleiveist2026ci-core-run,kleiveist2026ci-desktop-run,kleiveist2026acceptance}.
Daraus folgen vor allem Betriebs-, Plattform- und Prozessrisiken: Die
vorliegenden Paketergebnisse lassen sich nicht auf signierte Releases und
nicht auf einen produktiven Cloudbetrieb übertragen. Die Fallstudie bewertet
daher eine technische Template-Baseline, nicht die Produktionsreife eines
daraus entstehenden Produkts.
